% Generated by html2tex: Version 2.7 of February 6, 2012.
% Written by F.J. Faase.  http://www.iwriteiam.nl/

% html: Beginning of file: `test.html'
% DOCTYPE html
% [if IE 6]><![endif]
% [if IE 7]><![endif]

% <div id="globalWrapper">
		
% <div id="column-content">
			
% <div role="main">
												
% <div class="mw-indicators">

% </div class="mw-indicators">
				
\chapter{New features in this release}

\label{f0}								
% <div class="mw-body-content">

				
% <div class="intro">

% <span style="font-size:120%">
\textbf{This page is an overview of the key new functionality that has been introduced in Audacity 2.2.0}
% </span style="font-size:120%">

\begin{itemize}
\item  
% <span style="font-size:110%">
Details of all the major changes since 2.1.3 can be found in Release Notes 2.2.0.
% </span href="https://wiki.audacityteam.org/wiki/Release_Notes_2.2.0">

% </div href="https://wiki.audacityteam.org/wiki/Release_Notes_2.2.0">

\end{itemize}


\section{
% <span id="New_Logo">
New Logo
% </span id="New_Logo">
}


The logo has been given a refresh, and now uses a sans-serif font and a flatter style.


 \includegraphics[width=8.6cm]{../help/manual/m/images/8/88/audacity_logo_signika_512_transparent.png} 
 
Click to visit the Audacity main website for  the latest news and Audacity downloads\footnote{See URL http://audacityteam.org/}
% <div id="themes">

% </div id="themes">

\section{
% <span id="Themes">
Themes\protect\footnote{See URL http://alphamanual.audacityteam.org/man/Themes}
% </span title="Themes">
}


Audacity now comes supplied with four pre-configured, user-selectable, themes.  This enables you to choose the look and feel you prefer for Audacity's interface. see the Themes\footnote{See URL http://alphamanual.audacityteam.org/man/Themes} page for details.
\begin{itemize}
\item \textbf{Light} theme: this is a light theme loosely based on the look and feel of earlier Audacity versions, but given a contemporary twist with more modern-looking buttons and icons. 
\item \textbf{Dark} theme: created by the Dark Audacity project.\footnote{See URL http://www.darkaudacity.com/} This is similar to the Light theme, with the same buttons and icons, but given a dark twist.
\item \textbf{Classic} theme: The one you know and loved. This theme is a re-creation of the look and feel of earlier Audacity versions. 
\item \textbf{High Contrast} theme: some users with poor eyesight benefit from a high contrast that is 'eye-popping' for most people.
\end{itemize}

\* \* \* \* Theme Light.png\footnote{See URL http://alphamanual.audacityteam.org/man/Themes\#light}
\* \* \* \* Theme Dark.png\footnote{See URL http://alphamanual.audacityteam.org/man/Themes\#dark}
\* \* \* \* 
\textbf{Light} theme
\* \* \* \* 
\textbf{Dark} theme

\* \* \* \* Theme Classic.png\footnote{See URL http://alphamanual.audacityteam.org/man/Themes\#classic}
\* \* \* \* Theme Hi Contrast.png\footnote{See URL http://alphamanual.audacityteam.org/man/Themes\#hicontrast}
\* \* \* \* 
\textbf{Classic} theme
\* \* \* \* 
\textbf{High Contrast} theme


% <div class="note">
The theme to use can be chosen at 
% <span class="menu">
Edit \mbox{$>$} Preferences \mbox{$>$} Interface\footnote{See URL http://alphamanual.audacityteam.org/man/Interface.Preferences}
% </span title="Interface Preferences">
.  
\begin{itemize}
\item  In addition to the four pre-configured themes there is also a Custom theme.  
\end{itemize}

By default the custom theme looks the same as Classic theme - but, if you have the right programming skills and tools, you can use this template to create your own theme.  Instructions for how to do this may be found in the Audacity Wiki.
% </div href="https://wiki.audacityteam.org/wiki/Theming_HowTo">




% <div id="midi">

% </div id="midi">

\section{
% <span id="MIDI_.28and_Allegro.29_Playback">
MIDI (and Allegro) Playback
% </span id="MIDI_.28and_Allegro.29_Playback">
}


Playback of MIDI (and Allegro) files imported into Note Tracks\footnote{See URL http://alphamanual.audacityteam.org/man/Note.Tracks} is now available.  Please see the Note Tracks\footnote{See URL http://alphamanual.audacityteam.org/man/Note.Tracks} page for more details.

This should just work on Windows but for playback on Mac and Linux additional software may be required, see this section\footnote{See URL http://alphamanual.audacityteam.org/man/Playing.and.Recording\#midi} on the Playing and Recording\footnote{See URL http://alphamanual.audacityteam.org/man/Playing.and.Recording} page.

But note that there will no use of the Playback meter\footnote{See URL http://alphamanual.audacityteam.org/man/Meter.Toolbar\#playback} while Note tracks are played. 
% <div id="stemplots">

% </div id="stemplots">


\section{
% <span id="Stem_Plots">
Stem Plots
% </span id="Stem_Plots">
}


There is a new entry in the Tracks Preferences\footnote{See URL http://alphamanual.audacityteam.org/man/Tracks.Preferences} for \textbf{Display samples}.  This setting changes how Waveform\footnote{See URL http://alphamanual.audacityteam.org/man/Audacity.Waveform} and Waveform dB\footnote{See URL http://alphamanual.audacityteam.org/man/Audacity.Waveform\#db} views are displayed.  It only affects the appearance of the waveform when you are so far zoomed in that you can see the individual sample dots.  At lower zoom levels it makes no difference.  
\begin{itemize}
\item \textbf{Stem plot:} This is the default setting which draws a vertical line from the track center line to the sample dot, giving a clearer impression of the relative amplitude of the samples. As seen in the images below, when zoomed out close to the minimum for a stem plot, the horizontal distance between sample dots may be more uneven than seen with the connect dots default.   
\item \textbf{Connect dots:} This is alternative setting yields a waveform where each sample dot is connected to the next sample by a line drawn between them.  
\end{itemize}

% <span title="Examples of zoomed in Waveform view showing &#39;Connect dots&#39; and &#39;Stem Plot">
Connect dots \& Stem plot examples.png
% </span /="">
??

% <div class="note">
Uneven spacing is due to \texttt{{}"{}}aliasing\texttt{{}"{}}, but zooming in further will equalize the spacing whether choosing Connect dots or Stem plot.
% </div class="note">



% <div id="menus">

% </div id="menus">

\section{
% <span id="Menu_reorganization">
Menu reorganization
% </span id="Menu_reorganization">
}


We have made the Menus shorter and clearer than in previous Audacity versions. The menus have been simplified without losing functionality. The most commonly used functions are found in the top levels of the menus. The functions moved down into lower submenus are better organized.

This is not just a rearrangement.  We also added new menu items to make the layout more logical.  There are new menu items for exporting as MP3 or WAV. Previously you had to export audio, and then choose the format. You still can do that, but these new items are there for convenience.
% <div class="note">

One of the long standing bug-bears with Audacity is the distinction between \^aSave\^a and \^aExport\^a. People expect to be able to open a WAV audio file, edit it in Audacity and then click save. That is not how Audacity works. Audacity needs audio in its own unique format to work on it. So Audacity converts when you 'Open' a WAV file and converts back when you 'Save'. In an attempt to make this clearer Audacity uses the word 'Import' for opening a file like WAV and 'Export' for converting and saving in a format like WAV. Open and Save are reserved for Audacity's own project format.

Now the \^aExport\^a options are under the \^aSave Other\^a menu item, where people trying to save audio as an MP3 or WAV file are more likely to find them. 
% </div class="note">



% <div id="extendedmenubar">

% </div id="extendedmenubar">

\section{
% <span id="The_Extended_Menu_bar">
The Extended Menu bar
% </span id="The_Extended_Menu_bar">
}


There are two new additional menus that are hidden by default.  They can be turned on at 
% <span class="menu">
View \mbox{$>$} Extra Menus (on/off)\footnote{See URL http://alphamanual.audacityteam.org/man/View.Menu\#extra.menus.onoff}
% </span title="View Menu">
 or the Interface pane of 
% <span class="menu">
Edit \mbox{$>$} Preferences\footnote{See URL http://alphamanual.audacityteam.org/man/Interface.Preferences}
% </span title="Interface Preferences">
.

These extra menus have many extra less frequently used commands.  They are particularly useful to VI users, but normally-sighted users may find them useful too.

\textit{\textbf{Image of the Extended Menu bar as it appears on Windows}}
% <div class="noresize">

% <map name="ImageMap_1_1694869993">

% <area /="">

% <area /="">

% </map /="">
Menu Toolbar Extra.png
% </div /="">
??
\textit{Click, or hover, on either of the \textbf{Ext-} items at the end of the image to read about those menu entries.}

Menu
What you'll find there
Ext-Bar\footnote{See URL http://alphamanual.audacityteam.org/man/Ext-Bar.Menu}
The Ext-Bar menu provides access to Toolbar\footnote{See URL http://alphamanual.audacityteam.org/man/Toolbars.Overview} operations that are not available in the default Audacity menus. These will be of most interest to visually impaired users or those who have difficulty using the mouse.

Shortcuts\footnote{See URL http://alphamanual.audacityteam.org/man/Keyboard.Preferences} can be assigned to these commands if required.
Ext-Command\footnote{See URL http://alphamanual.audacityteam.org/man/Ext-Command.Menu}
The Ext-Command menu provides access to extra commands for track focus\footnote{See URL http://alphamanual.audacityteam.org/man/Audio.Tracks\#focus} and movement of the editing or playback cursor that are not available in the default Audacity menus. These will be of most interest to visually impaired users or those who have difficulty using the mouse.

Shortcuts\footnote{See URL http://alphamanual.audacityteam.org/man/Keyboard.Preferences} can be assigned to these commands if required.
% <div class="note">

If you only require regular access to a small set of these commands you can set shortcuts for them and leave the extra menus hidden.
% </div class="note">



% <div id="append">

% </div id="append">

\section{
% <span id="Appended_recording_on_the_same_track_in_now_the_default">
Appended recording on the same track in now the default
% </span id="Appended_recording_on_the_same_track_in_now_the_default">
}


From Audacity 2.2.0 onward the default recording mode has changed so that when you click the Record button The Record button on Transport Toolbar\footnote{See URL http://alphamanual.audacityteam.org/man/Transport.Toolbar}, or use the 
% <span style="background-color: #e0e0f0">
R
% </span style="background-color: #e0e0f0">
, Audacity will record at the end of the currently selected (or only) track.

\subsection{
% <span id="To_record_on_a_new_track">
To record on a new track
% </span id="To_record_on_a_new_track">
}


If you hold the 
% <span class="button">
Shift
% </span class="button">
 button down the Record button in Transport Toolbar\footnote{See URL http://alphamanual.audacityteam.org/man/Transport.Toolbar} will temporarily change to The Record New Track button. Then clicking on this modified Record button, or using the shortcut 
% <span style="background-color: #e0e0f0">
Shift + R
% </span style="background-color: #e0e0f0">
 will cause Audacity to create a new track and begin recording on that track from the current cursor position (or from the left edge of a region on the Timeline\footnote{See URL http://alphamanual.audacityteam.org/man/Timeline}). 


% <div id="help-buttons">

% </div id="help-buttons">

\section{
% <span id="Help_buttons">
Help buttons
% </span id="Help_buttons">
}


Many places in the user interface have had a help button Help Button\footnote{See URL http://alphamanual.audacityteam.org/man/Help.Button} added. Examples are all the Preferences dialog panes, all the Effects, Generators and Analyzers and some error messages.

Clicking on that button in the dialog will link you to the appropriate page in the Manual.


% <span style="font-weight: bold; font-size:118%">
Example:
% </span style="font-weight: bold; font-size:118%">
 the Amplify effect.  Try clicking on the \texttt{{}"{}}\textbf{?}\texttt{{}"{}} at the bottom right of this image.

% <div class="noresize">

% <map name="ImageMap_2_545976105">

% <area /="">

% </map /="">
Amplify.png
% </div /="">
??


% <div id="shortcuts">

% </div id="shortcuts">

\section{
% <span id="Standard_and_Full_shortcut_sets">
Standard and Full shortcut sets
% </span id="Standard_and_Full_shortcut_sets">
}


For Audacity 2.2.0 we have reduced the number of preset shortcuts in the application to a \texttt{{}"{}}\textbf{Standard\footnote{See URL http://alphamanual.audacityteam.org/man/Keyboard.Shortcut.Reference}}\texttt{{}"{}} set.  We did this to simplify the set of shortcuts somewhat and to provide greater flexibility for users who want to set their own custom shortcuts.

You can choose to revert to the full set of shortcuts that were in 2.1.3 and earlier by selecting \texttt{{}"{}}\textbf{Full}\texttt{{}"{}} from the dropdown menu accessed from the 
% <span class="button">
Defaults
% </span class="button">
 button in the Keyboard Preferences\footnote{See URL http://alphamanual.audacityteam.org/man/Keyboard.Preferences} dialog.

You can use the 
% <span class="button">
Defaults
% </span class="button">
 button to switch between the two provided default sets of shortcuts at any time.


% <div class="noresize">

% <map name="ImageMap_3_1572635641">

% <area /="">

% <area /="">

% </map /="">
Preferences Keyboard with Default dropdown menu - trimmed.png
% </div /="">
??

See Commands and Keyboard Shortcut Reference\footnote{See URL http://alphamanual.audacityteam.org/man/Keyboard.Shortcut.Reference} for more details.


% <div id="selection">

% </div id="selection">

\section{
% <span id="Selection_Toolbar_improvements">
Selection Toolbar improvements
% </span id="Selection_Toolbar_improvements">
}


There are now four available settings in the Selection and Audio Position Boxes in Selection Toolbar\footnote{See URL http://alphamanual.audacityteam.org/man/Selection.Toolbar} for the  manner in which the details of your selection are displayed:
\begin{itemize}
\item \textbf{Start and End of selection}: the start time and the end time of your selection \textit{(default setting)}
\item \textbf{Start and Length of selection}: the start time and the length of your selection
\item \textbf{Length and End of selection}: the length and the end time of your selection
\item \textbf{Length and Center of selection}: the length and the time at the center of your selection
\end{itemize}

% <span title="Selection Toolbar with dropdown menu for display type selection">
Selection Toolbar display modes.png
% </span /="">
??



% <div id="clips">

% </div id="clips">

\section{
% <span id="New_commands_for_using_clips_via_the_keyboard">
New commands for using clips via the keyboard
% </span id="New_commands_for_using_clips_via_the_keyboard">
}


New commands, all of which interact with the clips on the focused track. As yet, there are no default shortcuts:
\begin{enumerate}
\item 
% <span class="menu">
Transport \mbox{$>$} Cursor to \mbox{$>$} Previous clip boundary\footnote{See URL http://alphamanual.audacityteam.org/man/Transport.Menu:.Cursor.to\#previous.clip.boundary}
% </span title="Transport Menu: Cursor to">

\item 
% <span class="menu">
Transport Cursor to \mbox{$>$} Next clip boundary\footnote{See URL http://alphamanual.audacityteam.org/man/Transport.Menu:.Cursor.to\#next.clip.boundary}
% </span title="Transport Menu: Cursor to">

\item 
% <span class="menu">
Select \mbox{$>$} Clip Boundaries \mbox{$>$} Previous clip boundary to cursor\footnote{See URL http://alphamanual.audacityteam.org/man/Select.Menu:.Clip.Boundaries\#previous.clip.boundary.to.cursor} 
% </span title="Select Menu: Clip Boundaries">

\item 
% <span class="menu">
Select \mbox{$>$} Clip Boundaries \mbox{$>$} Cursor to next clip boundary\footnote{See URL http://alphamanual.audacityteam.org/man/Select.Menu:.Clip.Boundaries\#cursor.to.next.clip.boundary} 
% </span title="Select Menu: Clip Boundaries">

\item 
% <span class="menu">
Select \mbox{$>$} Clip Boundaries \mbox{$>$} Previous clip\footnote{See URL http://alphamanual.audacityteam.org/man/Select.Menu:.Clip.Boundaries\#previous.clip} 
% </span title="Select Menu: Clip Boundaries">

\item 
% <span class="menu">
Select \mbox{$>$} Clip Boundaries \mbox{$>$} Next clip\footnote{See URL http://alphamanual.audacityteam.org/man/Select.Menu:.Clip.Boundaries\#next.clip} 
% </span title="Select Menu: Clip Boundaries">

\end{enumerate}

% <div class="note">
There are also two new commands (Clip Left\footnote{See URL http://alphamanual.audacityteam.org/man/Ext-Command.Menu:.Cursor\#clip.left} and Clip Right\footnote{See URL http://alphamanual.audacityteam.org/man/Ext-Command.Menu:.Cursor\#clip.right}) for moving clips, which are available in the extended menus\footnote{See URL http://alphamanual.audacityteam.org/man/New.features.in.this.release\#extendedmenubar} added in this release.

Note that these extended menus are \textit{\textbf{not}} turned on by default, to turn them on please see Menu Reference\footnote{See URL http://alphamanual.audacityteam.org/man/Menu.Reference\#the.extended.menu.bar} for how to reveal the extended menus.
% </div title="Menu Reference">



% <div id="safety">

% </div id="safety">

\section{
% <span id="Running_out_of_disk_space">
Running out of disk space
% </span id="Running_out_of_disk_space">
}


We have now provided an error trap for situations where you are running out of available disk space.

You will now see the error message:
\texttt{{}"{}}\textbf{Audacity failed to write to a file in \mbox{$<$}device\mbox{$>$}}\texttt{{}"{}}

This is particularly useful when recording as Audacity will stop recording when the error is trapped, preserving your recording up to that point.


% <div id="appendix">

% </div id="appendix">

\section{
% <span id="Additional_new_features">
Additional new features\protect\footnote{See URL http://alphamanual.audacityteam.org/man/New.features.in.this.release.-.appendix}
% </span title="New features in this release - appendix">
}


This page \textbf{New features in this release - appendix\footnote{See URL http://alphamanual.audacityteam.org/man/New.features.in.this.release.-.appendix}} gives an overview of further new functionality that has been introduced in this release of Audacity.

